% HiMCM 2022 Problem A --- Non-technical grower / beekeeper letter (~1.5 pp)
% Compile: pdflatex grower_advisory_letter.tex
\documentclass[11pt]{article}
\usepackage[margin=1in]{geometry}
\usepackage{setspace}
\usepackage{enumitem}
\usepackage[hidelinks]{hyperref}
\onehalfspacing
\pagestyle{empty}

\begin{document}

\begin{center}
{\Large\bfseries Maximizing Orchard Yield and Apiary Sustainability:}\\[0.25em]
{\Large\bfseries A Data-Driven Guide for Commercial Pollination}\\[0.6em]
{\normalsize Prepared for commercial beekeeping associations and almond
growers}\\[0.2em]
{\small HiMCM 2022 Problem A \quad Team XXXX \quad (plain-language brief)}
\end{center}

\vspace{0.6em}
\noindent
Dear Association Members and Orchard Partners,

This short guide translates our colony and orchard model into field
practice. It is written for people who move bees and sign pollination
contracts---not for mathematicians. The numbers below come from our
simulation files for a \textbf{20-acre almond block}
(\(\approx 81{,}000\,\mathrm{m}^2\)).

\subsection*{1. How many hives? The golden band}
On a 20-acre almond orchard, plan on \textbf{1 to 2 hives per acre}, that
is \textbf{20 to 40 strong colonies} for the whole block. In our
optimization, \textbf{40 colonies} (2 per acre) gave the best balance of
nut set and rental cost at about \textbf{85\% of full pollination
potential} and roughly \textbf{\$51{,}000} net orchard return after a
\$200/hive rental. Pushing past about 2.5 hives/acre barely raises set
but starts to starve colonies of nectar and pollen (shared bloom), so
bees leave the contract weaker.

\begin{itemize}[leftmargin=1.2em]
  \item \textbf{Light bloom / small trees:} start near 1 hive/acre (20
        colonies).
  \item \textbf{Full bloom / mature block:} target 1.5--2 hives/acre
        (30--40 colonies).
  \item \textbf{Avoid:} stacking 3+ hives/acre as a default ``insurance''
        load without checking stores after petal fall.
\end{itemize}

\subsection*{2. Fan (sector) placement on the block}
Scatter colonies so flight paths cover the canopy instead of fighting
over the same fence line.

\begin{enumerate}[leftmargin=1.2em]
  \item Divide the 20 acres into \textbf{4--6 sectors} (roughly
        3--5 acres each).
  \item Place \textbf{groups of 4--8 pallets} near sector centers or
        along alleyways with truck access---not one giant wall of hives
        on a single edge.
  \item Face entrances \textbf{southeast} where wind and morning sun
        allow; keep \textbf{8--12 ft} between pallet groups when space
        allows.
  \item Keep a \textbf{water source} within a short flight; muddy
        puddles beat none at all during dry bloom weeks.
  \item Record GPS/pallet IDs so weak colonies can be swapped without
        reshuffling the whole yard.
\end{enumerate}

\subsection*{3. Feeding through bloom}
Almond bloom is early and can be cold. Even with orchard nectar, packages
often need help.

\begin{itemize}[leftmargin=1.2em]
  \item \textbf{Before move-in:} colonies should be queenright, with
        sealed brood and enough honey that frames are not light when
        tipped.
  \item \textbf{During cool bloom (\(<55^\circ\mathrm{F}\) flight days):}
        offer \textbf{1:1 sugar syrup} in hive-top feeders and a
        \textbf{pollen substitute patty} if pollen traps/frames look
        thin.
  \item \textbf{Stop heavy syrup} once consistent flight and orchard
        nectar flow are obvious---you want bees on flowers, not on
        feeders.
  \item \textbf{After petal fall:} inspect stores before the next
        contract; our model treats post-bloom nectar and brood survival
        as the hive-health score you feel as ``spring build-up.''
\end{itemize}

\subsection*{4. Varroa: the critical control that protects foragers}
Our stress tests show that when mites, spray exposure, and cold snaps
stack, colonies tip from a safe winter store zone into collapse as
\textbf{effective forager death} rises. You do not need the equations;
you need thresholds you can act on.

\begin{itemize}[leftmargin=1.2em]
  \item \textbf{Monitor:} alcohol wash or sticky board \textbf{before}
        almond staging and again after the contract.
  \item \textbf{Treat on the calendar that matches your region}, but do
        not enter pollination with high mite loads. As a practical
        trigger used by many commercial yards: act when washes approach
        the \textbf{high single digits per 300 bees} (follow your
        association's current economic threshold---treat \emph{before}
        the contract if you are already close).
  \item \textbf{Rotate modes of action}; never rely on a single chemical
        year after year.
  \item \textbf{Coordinate with the grower} on fungicide/insecticide
        timing. Night sprays and avoiding open-bloom insecticides
        protect the foragers your rent is paying for.
\end{itemize}

\subsection*{5. Shared checklist (grower + beekeeper)}
\begin{enumerate}[leftmargin=1.2em]
  \item Contract states hive count in the \textbf{20--40} band for 20
        acres, colony strength standard, and payment timing.
  \item Map shows \textbf{sector/fan placement} and water.
  \item Spray schedule shared \textbf{72 hours} ahead when possible.
  \item Post-bloom: store check + mite check before the next move.
\end{enumerate}

\vspace{0.4em}
\noindent
Healthy bees and full almond sets are the same project. Stock the golden
band, place colonies so they cover the canopy, feed only when weather
steals flight, and keep Varroa from quietly raising forager death during
the most valuable weeks of the year.

\vspace{0.8em}
\noindent
Respectfully,\\[0.6em]
Team XXXX\\
HiMCM 2022 Problem A Modeling Group

\end{document}
